\documentclass[11pt,letterpaper]{article}

\usepackage[top=1in, bottom=1in, left=1in, right=1in, headheight=14pt]{geometry}
\usepackage{fontspec}
\setmainfont{DejaVu Serif}
\setsansfont{DejaVu Sans}
\setmonofont{DejaVu Sans Mono}

\usepackage{xcolor}
\usepackage{titlesec}
\usepackage{fancyhdr}
\usepackage{enumitem}
\usepackage{hyperref}
\usepackage{booktabs}
\usepackage{tabularx}
\usepackage{longtable}
\usepackage{colortbl}
\usepackage{multirow}
\usepackage{array}
\usepackage{parskip}
\usepackage{tcolorbox}
\usepackage{graphicx}
\usepackage{lastpage}

%% ── COLOR SCHEME: Technology ───────────────────────────────────────────────
\definecolor{primary}{HTML}{263238}
\definecolor{secondary}{HTML}{1565C0}
\definecolor{tertiary}{HTML}{37474F}
\definecolor{accent}{HTML}{1976D2}
\definecolor{lightprimary}{HTML}{ECEFF1}
\definecolor{lightsecondary}{HTML}{E3F2FD}
\definecolor{lighttertiary}{HTML}{ECEFF1}
\definecolor{rowgray}{HTML}{F5F5F5}
\definecolor{bordergray}{HTML}{BDBDBD}
\definecolor{subtlegray}{HTML}{757575}
\definecolor{darktext}{HTML}{212121}
\definecolor{warnred}{HTML}{B71C1C}
\definecolor{okgreen}{HTML}{1B5E20}

%% ── SECTION FORMAT ──────────────────────────────────────────────────────────
\titleformat{\section}
  {\Large\bfseries\sffamily\color{primary}}
  {\thesection}{1em}{}
  [\vspace{-0.5em}\textcolor{bordergray}{\rule{\textwidth}{0.4pt}}]

\titleformat{\subsection}
  {\large\bfseries\sffamily\color{secondary}}
  {\thesubsection}{1em}{}

\titleformat{\subsubsection}
  {\normalsize\bfseries\sffamily\color{tertiary}}
  {\thesubsubsection}{1em}{}

\titlespacing*{\section}{0pt}{1.5em}{0.8em}
\titlespacing*{\subsection}{0pt}{1.2em}{0.5em}
\titlespacing*{\subsubsection}{0pt}{0.8em}{0.3em}

%% ── CUSTOM ENVIRONMENTS ─────────────────────────────────────────────────────
\newcommand{\stageheader}[2]{%
  \noindent\colorbox{#2}{\makebox[\textwidth][l]{%
    \textcolor{white}{\bfseries\sffamily\hspace{6pt}#1\hspace{6pt}}}}%
  \vspace{0.5em}\par%
}

\tcbuselibrary{breakable,skins}

\newtcolorbox{asciibox}{
  colback=rowgray, colframe=bordergray,
  arc=1pt, boxrule=0.5pt,
  left=6pt, right=6pt, top=4pt, bottom=4pt,
  fontupper=\small\ttfamily,
  breakable
}

\newtcolorbox{quotebox}{
  colback=lightprimary, colframe=primary,
  arc=2pt, boxrule=0.5pt,
  left=12pt, right=12pt, top=8pt, bottom=8pt,
  breakable
}

\newcommand{\metafield}[2]{\textbf{\sffamily #1} #2\par}
\newcolumntype{L}[1]{>{\raggedright\arraybackslash}p{#1}}
\newcolumntype{C}[1]{>{\centering\arraybackslash}p{#1}}
\newcommand{\tableheader}[1]{\textbf{\sffamily\textcolor{white}{#1}}}

%% ── HEADERS / FOOTERS ───────────────────────────────────────────────────────
\pagestyle{fancy}
\fancyhf{}
\fancyhead[L]{\small\sffamily\textcolor{subtlegray}{Technical Documentation: Policy \& Procedures}}
\fancyhead[R]{\small\sffamily\textcolor{subtlegray}{GSD Mission Package}}
\fancyfoot[C]{\small\sffamily\textcolor{subtlegray}{Page \thepage\ of \pageref{LastPage}}}
\renewcommand{\headrulewidth}{0.4pt}
\renewcommand{\footrulewidth}{0pt}

%% ── HYPERLINKS ──────────────────────────────────────────────────────────────
\hypersetup{
  colorlinks=true,
  linkcolor=secondary,
  urlcolor=secondary,
  pdfauthor={GSD Ecosystem / Tibsfox Inc.},
  pdftitle={Technical Documentation: Policy \& Procedures -- Mission Package},
  pdfsubject={Vision to Mission Pipeline},
}

%% ════════════════════════════════════════════════════════════════════════════
\begin{document}

%% ── TITLE PAGE ──────────────────────────────────────────────────────────────
\thispagestyle{empty}
\begin{center}
\vspace*{3cm}

{\small\sffamily\textcolor{primary}{\textbf{GSD RESEARCH MISSION PACKAGE}}}

\vspace{0.3cm}
\textcolor{bordergray}{\rule{8cm}{0.4pt}}

\vspace{1.5cm}
{\fontsize{26}{32}\selectfont\bfseries\sffamily\textcolor{primary}{Technical Documentation\\[0.3em]for Policy \& Procedures}}

\vspace{0.8cm}
{\Large\sffamily\textcolor{secondary}{Engineering, Creating, and Governing\\Technical Documentation Systems}}

\vspace{0.5cm}
{\large\sffamily\itshape\textcolor{tertiary}{A Deep Research Study}}

\vspace{3cm}
{\sffamily\textcolor{accent}{Vision $\rightarrow$ Research $\rightarrow$ Mission}}\\[0.3em]
{\small\sffamily\textcolor{accent}{Full Pipeline Execution}}

\vspace{3cm}
{\small\sffamily\textcolor{subtlegray}{Date: April 5, 2026 \quad|\quad Status: Ready for GSD Orchestrator}}\\[0.3em]
{\small\sffamily\textcolor{subtlegray}{Activation Profile: Squadron (12 roles) \quad|\quad Estimated: 5--8 context windows across 2--3 sessions}}
\end{center}
\newpage

%% ── TABLE OF CONTENTS ───────────────────────────────────────────────────────
\tableofcontents
\newpage

%% ════════════════════════════════════════════════════════════════════════════
%%  STAGE 1 — VISION DOCUMENT
%% ════════════════════════════════════════════════════════════════════════════
\stageheader{STAGE 1 --- VISION DOCUMENT}{primary}

\section{Technical Documentation for Policy \& Procedures --- Vision Guide}

\metafield{Date:}{April 5, 2026}
\metafield{Status:}{Pre-Execution --- Ready for GSD Orchestrator}
\metafield{Depends On:}{gsd-skill-creator v1.49.x dev branch}
\metafield{Context:}{Cold-start research mission. No prior conversation on this specific topic.}
\metafield{Author:}{Miles Tiberius Foxglove / Tibsfox, Inc.}

\subsection{Vision}

Documentation is the memory of a system. When it is absent, degraded, or disconnected from the reality of the living codebase or operating procedure, every engineer who follows pays the debt. The Amiga shipped with a complete, lovingly-written developer reference because its creators understood that architecture is nothing without the map that helps others navigate it. The same principle applies to every policy, every procedure, every engineering standard ever written: the artifact that captures the \emph{why} and the \emph{how} is as load-bearing as the system it describes.

This mission addresses a persistent gap in the GSD ecosystem and in engineering organizations broadly: the discipline of creating, governing, and maintaining technical documentation for policies and procedures. Not documentation as afterthought --- documentation as first-class engineering deliverable. The shift is architectural: from ad-hoc document dumps to structured, version-controlled, AI-augmented documentation pipelines that travel with the system they describe.

The hierarchy matters. A policy is a management-intent statement. A standard is a quantifiable requirement. A procedure is the step-by-step execution. Confusing these three is not a semantic quibble --- it produces documentation that neither auditors nor engineers can act on. Getting the taxonomy right is the foundation on which every other best practice rests.

The current moment is a forcing function. With 64\% of software development professionals now using AI for documentation (Google Cloud DORA 2025), the question is no longer whether AI belongs in the documentation pipeline --- it is how to ensure that AI-generated content does not drift from truth. Docs-as-code, CI/CD pipeline integration, drift detection, and human-in-the-loop review are the answers the field has converged on. This mission maps the landscape, assembles the evidence, and packages it for GSD execution.

\subsection{Problem Statement}

\begin{enumerate}[leftmargin=*, label=\textbf{P\arabic*.}]

\item \textbf{Documentation Decay:} Technical documentation is created once and rarely updated, leading to progressively wider gaps between what the docs say and what the system does. Engineers stop trusting docs; workarounds proliferate; institutional knowledge migrates into individual heads and disappears when those people leave.

\item \textbf{Taxonomy Confusion:} The terms ``policy,'' ``standard,'' ``procedure,'' and ``guideline'' are routinely abused as synonyms across engineering organizations. This prevents compliance audits from finding what they need, prevents engineers from knowing what is binding versus advisory, and produces documentation structures that satisfy no one.

\item \textbf{Missing Governance Pipeline:} Documentation lacks the CI/CD discipline applied to code. There is no lint, no test, no merge gate, no drift detection. A changed API or altered procedure has no automated signal to the documentation system. The result is that documentation and reality diverge silently.

\item \textbf{Audience Mismatch:} Technical documentation written for one audience (e.g., senior engineers) is distributed to another (e.g., new hires, operations staff, auditors). Without persona-driven segmentation and modular topic design, single monolithic documents fail all of their audiences simultaneously.

\item \textbf{AI Without Guardrails:} AI-assisted documentation generation is now widespread but frequently misapplied --- generating ``correct but useless'' descriptions of what code does without capturing architectural intent, edge case rationale, or the \emph{why} behind design decisions. Human review gates are absent, and documentation drift goes undetected until an incident surfaces it.

\end{enumerate}

\subsection{Core Concept}

\textbf{Documentation as a Living System:} Every document is a versioned artifact with an owner, an audience, a lifecycle state, and a connection to the code or procedure it describes. Documentation does not exist apart from the system --- it is built, tested, reviewed, and deployed alongside it.

\subsubsection{Documentation Ecosystem Map}

\begin{asciibox}
DOCUMENTATION HIERARCHY (NIST / ISO / ISACA aligned)
=====================================================

  [Board / CEO]
       |
   POLICY  ─────── High-level management intent; strategic
       |            "What we will do and why"
   STANDARD ─────── Quantifiable requirements; specific controls
       |            "Exactly how much / how often / what threshold"
   PROCEDURE ─────── Step-by-step execution; operational
       |            "Do this, then this, then verify this"
   GUIDELINE ─────── Advisory; non-binding best practice
                    "Recommended approach; judgment allowed"

DOCUMENT LIFECYCLE
==================

  DRAFT ──► REVIEW ──► APPROVED ──► PUBLISHED ──► MAINTAINED
     ▲                                                  |
     └──────────────── UPDATED / RETIRED ◄──────────────┘

DOCS-AS-CODE PIPELINE
======================

  Author (Markdown/XML/DITA)
       |
  Git Branch ──► Pull Request ──► Lint / Style Check
                                        |
                                  Reviewer Approval
                                        |
                               Merge to main branch
                                        |
                           CI/CD Build + Publish
                                        |
                          Drift Detection (90-day check)
                                        |
                            Stale Flag ──► Owner Notified
\end{asciibox}

\subsubsection{Research Module Architecture}

\begin{asciibox}
MODULE DEPENDENCY GRAPH
========================

  M1: FOUNDATIONS      M2: STRUCTURED       M3: DOCS-AS-CODE
  Taxonomy, lifecycle  AUTHORING            Pipeline, CI/CD,
  governance, policy   DITA, OASIS,         drift detection,
  hierarchy            single-source        version control
       |    \_______________/                     |
       |          |                               |
       v          v                               v
  M4: COMPLIANCE FRAMEWORKS         M5: AI-AUGMENTED WRITING
  NIST, ISO, FISMA, CMMC            Generation, review gates,
  policy-to-control mapping         co-pilot best practices

  Cross-module synthesis: Governance model that unifies
  M1 taxonomy + M3 pipeline + M4 compliance + M5 AI guardrails
\end{asciibox}

\subsection{Research Modules}

\subsubsection{Module 1: Foundations \& Taxonomy}
Establishes the authoritative hierarchy of documentation types: policy, standard, procedure, guideline, control. Maps the document lifecycle from draft through retirement. Covers governance structures, ownership models, audience segmentation by persona, and the linguistic rules (active voice, second person for procedures, third person for formal documents) that make technical documentation actionable.

\subsubsection{Module 2: Structured Authoring Systems}
Surveys the DITA (Darwin Information Typing Architecture) standard and the broader structured authoring landscape. Covers topic-based authoring, content reuse via conref/conkeyref, single-source multi-format publishing, modular documentation design, and the tradeoffs between XML-based DITA and lighter-weight Markdown-based approaches. Evaluates when each approach is appropriate.

\subsubsection{Module 3: Docs-as-Code Pipeline}
Maps the engineering discipline of treating documentation like code: Git version control, branch-per-feature workflows, pull request review, CI/CD pipeline integration, automated linting and style checking, link validation, and drift detection. Covers tooling (GitBook, Mintlify, MkDocs, Hugo, Sphinx) and the metrics that make documentation quality measurable.

\subsubsection{Module 4: Compliance \& Regulatory Frameworks}
Surveys the major frameworks that govern technical documentation for policy and procedures: NIST SP 800-53, NIST CSF 2.0, NIST SP 800-100, ISO 27001/27002, FISMA, CMMC/NIST 800-171. Maps the specific documentation artifacts each framework requires (System Security Plans, POA\&Ms, Policies, Standards, Procedures) and how they interrelate. Covers the distinction between mandatory compliance requirements (MCR) and desired security requirements (DSR).

\subsubsection{Module 5: AI-Augmented Documentation}
Assesses the current state of AI in the documentation pipeline: generation (GitHub Copilot, IBM watsonx, Swimm, Apidog), review gates, drift detection via CI/CD hooks, and the human-in-the-loop patterns that keep AI-generated content accurate and useful. Identifies failure modes (``correct but useless'' generation, missing architectural intent) and the best-practice mitigations. References Google DORA 2025 data: 64\% of development professionals now use AI for documentation.

\subsection{Chipset Configuration}

\begin{asciibox}
GSD CHIPSET — TECHNICAL DOCUMENTATION MISSION
===============================================

  model: claude-opus-4-5     # Synthesis, governance design,
  role: Agnus (DMA)          # compliance mapping, cross-module
  allocation: ~30%           # integration, judgment-heavy analysis

  model: claude-sonnet-4-5   # Structured content, module surveys,
  role: Denise (Display)     # document assembly, verification,
  allocation: ~60%           # DITA reference, pipeline specs

  model: claude-haiku-3-5    # Shared schemas, templates,
  role: Paula (Audio)        # scaffolding, taxonomy tables
  allocation: ~10%           # wave 0 foundation work

  activation_profile: Squadron
  parallel_tracks: 2
  capcom_gates: [after Wave 0, after Wave 1, before publication]
\end{asciibox}

\subsection{Scope Boundaries}

\subsubsection{In Scope (v1.0)}
\begin{itemize}[leftmargin=*]
  \item Policy/standard/procedure/guideline taxonomy and hierarchy models
  \item DITA structured authoring architecture and topic-based design
  \item Docs-as-code pipeline implementation patterns (Git, CI/CD, linting)
  \item Major compliance framework documentation requirements (NIST, ISO, CMMC)
  \item AI-augmented documentation tools, patterns, and failure modes
  \item Document lifecycle governance: ownership, review cadence, drift detection
  \item Audience persona design and content segmentation strategies
  \item Version control patterns specific to documentation (not code)
  \item GSD-compatible output: skill templates for documentation generation
\end{itemize}

\subsubsection{Out of Scope}
\begin{itemize}[leftmargin=*]
  \item Legal document creation (contracts, IP filings) --- separate mission
  \item Specific SOP content for particular industries (healthcare SOPs, aviation checklists)
  \item Full compliance audit methodology --- requires compliance-specific mission
  \item User-facing documentation UX design (information architecture beyond structure)
  \item Localization and translation workflows
\end{itemize}

\subsection{Success Criteria}

\begin{enumerate}[leftmargin=*, label=\textbf{SC\arabic*.}]
  \item \textbf{Taxonomy Clarity:} Output includes a complete, authoritative documentation hierarchy (policy/standard/procedure/guideline/control) with definitions sourced from NIST, ISO, and ISACA.
  \item \textbf{DITA Coverage:} Module 2 produces a complete reference to DITA topic types (concept, task, reference, glossary) with decision criteria for when to use DITA vs. lighter Markdown-based approaches.
  \item \textbf{Pipeline Specification:} Module 3 produces a complete docs-as-code pipeline spec with concrete tooling recommendations, CI/CD job templates, and drift detection thresholds.
  \item \textbf{Framework Mapping:} Module 4 produces a cross-reference table mapping NIST 800-53 control families to required documentation artifacts, usable as an audit checklist.
  \item \textbf{AI Pattern Library:} Module 5 produces a pattern library of AI documentation use cases with associated guardrail requirements and review gate specifications.
  \item \textbf{GSD Skill Output:} Mission produces at least one GSD-compatible skill template for documentation generation (policy, procedure, or SOP template with structured fields).
  \item \textbf{Governance Model:} Synthesis produces a complete document governance model covering ownership, review cadence, lifecycle states, and escalation paths.
  \item \textbf{Persona Coverage:} At least three distinct audience personas are defined (e.g., senior engineer, new hire, compliance auditor) with content segmentation guidance per persona.
  \item \textbf{Source Quality:} All numerical claims are attributed to government agencies, peer-reviewed research, or professional organizations. Zero entertainment media or unsourced claims.
  \item \textbf{Compilation Clean:} LaTeX compiles in three passes without errors. PDF is complete, navigable, and includes working table of contents.
  \item \textbf{Test Coverage:} Test plan includes at least 2 tests per success criterion, with safety-critical tests marked mandatory-pass.
  \item \textbf{Handoff Ready:} Package is self-contained --- any GSD operator can execute the mission with zero additional context.
\end{enumerate}

\subsection{Relationship to Other Documents}

\begin{center}
\rowcolors{2}{rowgray}{white}
\begin{tabularx}{\textwidth}{L{4cm}L{4cm}X}
\toprule
\rowcolor{primary}
\tableheader{Document} & \tableheader{Relationship} & \tableheader{Notes} \\
\midrule
GSD Skill Creator v1.49.x & Parent system & Documentation skills will be added as GSD skills for orchestrator use \\
The Space Between & Philosophical spine & Mathematical foundations; autodidact methodology informs documentation pedagogy \\
GSD Vision-to-Mission Skill & Upstream dependency & This mission was generated via the vision-to-mission pipeline \\
NASA Mission Series & Parallel track & Documentation standards for mission packages are a direct application of this research \\
Fox Infrastructure Group & Downstream consumer & FIG operational SOPs will be engineered using outputs from this mission \\
College of Knowledge (.college/) & Downstream consumer & Rosetta Core documentation patterns depend on structured authoring research \\
\bottomrule
\end{tabularx}
\end{center}

\subsection{The Through-Line}

The Amiga's chipset didn't just run code faster --- it offloaded specialized tasks to dedicated silicon, freeing the CPU to do what it was best at. Technical documentation works the same way when it is built correctly: a structured, versioned, governed documentation system offloads cognitive load from the engineers who need to act on it, freeing them to solve actual problems rather than excavating institutional memory.

This is the ``giving people their lives back'' principle applied to knowledge infrastructure. Every hour a new engineer spends chasing an undocumented procedure, every compliance audit that fails because the policy hierarchy is ambiguous, every incident that recurs because the post-mortem was never properly documented --- these are the costs of treating documentation as overhead. This mission treats it as architecture. The spaces between the code are where the meaning lives.

%% ════════════════════════════════════════════════════════════════════════════
%%  STAGE 2 — RESEARCH REFERENCE
%% ════════════════════════════════════════════════════════════════════════════
\newpage
\stageheader{STAGE 2 --- RESEARCH REFERENCE}{secondary}

\section{Technical Documentation for Policy \& Procedures --- Research Reference}

\subsection{How to Use This Document}

This reference assembles findings from professional organizations, government agencies, and peer-reviewed sources. Each module section corresponds to a parallel research track in Stage 3. Agents assigned to a module should read the corresponding section before beginning work.

\textbf{Key source organizations:}
\begin{itemize}[leftmargin=*]
  \item \textbf{NIST} (National Institute of Standards and Technology) --- SP 800-53, SP 800-100, CSF 2.0
  \item \textbf{OASIS} --- DITA standard (Darwin Information Typing Architecture, current v1.3)
  \item \textbf{ISO/IEC} --- ISO 27001/27002 information security management
  \item \textbf{ISACA} --- Control objectives, governance terminology
  \item \textbf{Google Cloud DORA Research} --- Annual developer productivity and documentation surveys
  \item \textbf{Stack Overflow} --- Annual Developer Survey (2025 edition)
  \item \textbf{USU Engineering Writing Center} --- Technical writing standards
\end{itemize}

\subsection{Module 1 Reference: Foundations \& Taxonomy}

\subsubsection{The Documentation Hierarchy}

Industry-recognized terminology from NIST, ISO, ISACA, and AICPA establishes a clear hierarchy. These terms are \emph{not} synonymous, and using them interchangeably is a documented failure mode in compliance and engineering environments.

\begin{center}
\rowcolors{2}{rowgray}{white}
\begin{tabularx}{\textwidth}{L{2.5cm}L{4cm}L{3cm}X}
\toprule
\rowcolor{primary}
\tableheader{Term} & \tableheader{Definition} & \tableheader{Origin} & \tableheader{Example} \\
\midrule
Policy & High-level statement of management intent; formally establishes requirements & NIST 800-53, ISO 27001 & ``All production systems must be patched within 30 days of a critical CVE'' \\
Standard & Formally-established, quantifiable requirement; objective and measurable & NIST, ISACA & ``Passwords must be minimum 12 characters, rotated every 90 days'' \\
Procedure & Step-by-step operational instructions for executing a standard & NIST 800-100 & ``To patch a production server: 1. Open change ticket. 2. Schedule maintenance window...'' \\
Guideline & Advisory, non-binding recommended approach; judgment allowed & ISO 27002 & ``Consider using ephemeral credentials for CI/CD pipelines'' \\
Control & Any process, policy, device, or practice that modifies risk & ISO 27005, NIST & MFA implementation, firewall rule, access control list \\
\bottomrule
\end{tabularx}
\end{center}

\subsubsection{Document Lifecycle States}

Effective governance requires explicit lifecycle state tracking. Documents without defined states accumulate silently in ``unknown'' status, which is functionally the same as no governance at all.

\begin{center}
\rowcolors{2}{rowgray}{white}
\begin{tabularx}{\textwidth}{L{3cm}XL{3cm}}
\toprule
\rowcolor{primary}
\tableheader{State} & \tableheader{Definition} & \tableheader{Gate to Next} \\
\midrule
Draft & Under initial authorship; not for distribution & Owner review \\
Review & Circulated to stakeholders; open for comment & SME + compliance sign-off \\
Approved & Formally accepted by document owner and governance body & Publication date set \\
Published & Active, authoritative, accessible to audience & Scheduled review cycle \\
Under Review & Triggered by change event or scheduled review & Update or retire decision \\
Retired & Superseded or obsolete; archived, not deleted & Archive timestamp \\
\bottomrule
\end{tabularx}
\end{center}

\subsubsection{Audience Personas and Content Segmentation}

Good technical documentation is persona-driven. The Home Office Engineering Guidance and USU Engineering Writing Center both emphasize: write down who your users are and how technical they are before writing a single word.

\begin{center}
\rowcolors{2}{rowgray}{white}
\begin{tabularx}{\textwidth}{L{3cm}L{3cm}L{3cm}X}
\toprule
\rowcolor{primary}
\tableheader{Persona} & \tableheader{Goal} & \tableheader{Reading Mode} & \tableheader{Content Needs} \\
\midrule
Senior Engineer & Understand design rationale & Scan + deep read & Architecture decisions, constraints, edge cases \\
New Hire & Get operational quickly & Step-by-step & SOPs, runbooks, setup guides, glossary \\
Compliance Auditor & Verify controls are documented & Search by control ID & Policy hierarchy, control mappings, evidence artifacts \\
Operations Staff & Execute repeatable tasks & Follow procedure & Step-by-step procedures, checklists, decision trees \\
Product Manager & Understand system capabilities & Overview + summary & High-level summaries, capability maps \\
\bottomrule
\end{tabularx}
\end{center}

\subsubsection{Technical Writing Standards}

From USU Engineering Writing Center and UK Home Office Engineering Guidance:
\begin{itemize}[leftmargin=*]
  \item Use \textbf{active voice} for clarity and concision: ``The technician cleans the contacts'' not ``The contacts are cleaned by the technician''
  \item Use \textbf{second person} (``you'') for instructions and procedures
  \item Use \textbf{third person} for formal documents (proposals, design reports)
  \item Use \textbf{plain English}: complexity in the subject does not require complexity in the prose
  \item Apply verb tense consistently; shift tense only to reflect genuine changes in time
  \item Avoid clichés, slang, and antiquated technical terminology
  \item Each document entry should address a \textbf{single task, topic, or question} --- modular and focused
\end{itemize}

\subsection{Module 2 Reference: Structured Authoring Systems}

\subsubsection{DITA (Darwin Information Typing Architecture)}

DITA is an open XML-based standard maintained by OASIS (current version 1.3, approved December 2015). It defines topic-oriented, information-typed content that can be reused and single-sourced.

\textbf{Core DITA topic types:}
\begin{center}
\rowcolors{2}{rowgray}{white}
\begin{tabularx}{\textwidth}{L{3cm}XL{3.5cm}}
\toprule
\rowcolor{primary}
\tableheader{Topic Type} & \tableheader{Purpose} & \tableheader{Use Case} \\
\midrule
Concept & Background information; understanding-oriented & System overviews, theory sections \\
Task & Step-by-step procedures; action-oriented & SOPs, how-to guides, runbooks \\
Reference & Look-up information; data tables, specifications & API references, parameter tables \\
Troubleshooting & Symptom-cause-remedy structure & Known issue resolution \\
Glossary Entry & Single-term definitions & Terminology sections \\
\bottomrule
\end{tabularx}
\end{center}

\textbf{Key DITA mechanisms:}
\begin{itemize}[leftmargin=*]
  \item \textbf{Conref / Conkeyref:} Transclusion --- reference content from one topic into another; update once, reflects everywhere
  \item \textbf{DITA Maps:} Container for topics; defines sequence, structure, hierarchy, and navigation
  \item \textbf{Conditional Processing (.ditaval):} Filter content by audience, platform, or product without duplicating topics
  \item \textbf{Specialization:} Extend base DITA types for domain-specific needs while retaining standard processing
  \item \textbf{Relationship Tables (reltables):} Define hyperlinks between topics; maintain cross-references automatically
\end{itemize}

\textbf{DITA vs. Markdown decision criteria:}

\begin{center}
\rowcolors{2}{rowgray}{white}
\begin{tabularx}{\textwidth}{L{3cm}L{4cm}L{4cm}}
\toprule
\rowcolor{primary}
\tableheader{Factor} & \tableheader{Use DITA} & \tableheader{Use Markdown / Docs-as-Code} \\
\midrule
Scale & Enterprise, 1000s of topics & Small-to-medium teams \\
Reuse needs & High (single-source, multi-output) & Low-to-medium \\
Audience & Product documentation, regulated industries & Developer docs, wikis, runbooks \\
Tooling budget & High (oXygen, Arbortext, etc.) & Low (any text editor + Git) \\
Translation & Required (XLIFF integration) & Limited tooling \\
Compliance & Strict (ISO, FDA, aerospace) & Flexible environments \\
\bottomrule
\end{tabularx}
\end{center}

\subsubsection{Content Reuse Patterns}

The LEGO blocks principle: create standard components (product overviews, warning messages, common procedures) and assemble them into different guides as needed. When a step changes, update once --- all publications reflect the change automatically. This is architecturally equivalent to DRY (Don't Repeat Yourself) in code.

\subsection{Module 3 Reference: Docs-as-Code Pipeline}

\subsubsection{Core Principle}

Docs-as-code applies software development discipline to documentation: write in plain text (Markdown, reStructuredText, or DITA XML), version in Git, review via pull requests, deploy via CI/CD pipelines, and test for quality automatically.

\textbf{Business impact:} Well-written documentation can reduce support tickets by 40--60\% (Paligo, 2025). IBM internal testing found AI-assisted documentation reduced documentation time by an average of 59\%.

\subsubsection{Pipeline Architecture}

\begin{asciibox}
DOCS-AS-CODE CI/CD PIPELINE
=============================

Author writes Markdown / XML
         |
  git checkout -b docs/feature-name
  [edit] → git commit -m "docs: add failover procedure"
         |
  git push → Pull Request opened
         |
  CI PIPELINE TRIGGERS:
  ├── Lint (Vale, markdownlint)       ← style guide enforcement
  ├── Link validation                 ← no broken cross-references
  ├── Spell check                     ← error detection
  ├── Build (Hugo / MkDocs / DITA-OT) ← verify publishable
  └── Drift check                     ← docs older than threshold?
         |
  Reviewer approval (PR review)
         |
  Merge to main → Auto-deploy to docs site
         |
  Scheduled drift detection:
  └── Flag docs >90 days old in modules with recent commits
      └── Owner notified → update or retire decision
\end{asciibox}

\subsubsection{Tooling Landscape}

\begin{center}
\rowcolors{2}{rowgray}{white}
\begin{tabularx}{\textwidth}{L{3cm}L{3cm}XL{2.5cm}}
\toprule
\rowcolor{primary}
\tableheader{Tool} & \tableheader{Category} & \tableheader{Strengths} & \tableheader{Best For} \\
\midrule
GitBook & Publishing platform & Git-native, SOC 2 + ISO 27001, CI/CD integration & Developer portals, API docs \\
MkDocs + Material & Static site generator & Python-native, fast, excellent search & Engineering wikis, runbooks \\
Hugo & Static site generator & Fastest build times, flexible & Large doc sites \\
Sphinx & Static site generator & Deep Python ecosystem integration & Open source projects \\
Mintlify & AI-augmented platform & Auto-update from code changes & API documentation \\
Vale & Linter & Style guide enforcement (Google, Microsoft style guides built in) & Any Markdown/RST repo \\
Swimm & AI drift detection & Attaches docs to code logic at commit time & Legacy codebase documentation \\
DITA-OT & DITA processor & Open-source, multi-format output (PDF, HTML, EPUB) & Enterprise DITA deployments \\
\bottomrule
\end{tabularx}
\end{center}

\subsubsection{Drift Detection Metrics}

Documentation drift is the silent killer of documentation systems. From the 2026 AI-Driven Documentation analysis (Overcast Blog):
\begin{itemize}[leftmargin=*]
  \item Configure CI pipeline to report: ``\texttt{x\% of changed files have no corresponding doc update}''
  \item Flag docs older than a threshold (e.g., 90 days) in modules with recent commits
  \item Block merges when drift risk exceeds threshold, just as test failures block deploys
  \item Use PR hooks: when a PR touches a method already linked in the documentation map, prompt ``Does the associated doc need updating?''
\end{itemize}

\subsection{Module 4 Reference: Compliance \& Regulatory Frameworks}

\subsubsection{Framework Landscape}

\begin{center}
\rowcolors{2}{rowgray}{white}
\begin{tabularx}{\textwidth}{L{3cm}L{3cm}XL{2.5cm}}
\toprule
\rowcolor{primary}
\tableheader{Framework} & \tableheader{Owner} & \tableheader{Scope} & \tableheader{Best For} \\
\midrule
NIST CSF 2.0 & NIST & Risk-based cybersecurity guidance; least prescriptive & Small organizations, voluntary alignment \\
NIST SP 800-53 R5 & NIST & Comprehensive federal security control catalog & Federal agencies, DoD contractors \\
NIST SP 800-171 / CMMC & NIST / DoD & Protecting Controlled Unclassified Information (CUI) & Defense Industrial Base \\
ISO 27001/27002 & ISO/IEC & International ISMS certification & International organizations \\
NIST SP 800-100 & NIST & Information security governance handbook & Agency governance programs \\
FISMA & US Congress & Federal information system security requirements & Federal agencies \\
\bottomrule
\end{tabularx}
\end{center}

\subsubsection{Required Documentation Artifacts by Framework}

NIST SP 800-53 requires specific policy and procedure documents for each of its 20 control families. Every control family requires:
\begin{itemize}[leftmargin=*]
  \item A \textbf{policy} document: high-level management intent statement
  \item A \textbf{procedures} document: step-by-step implementation of the policy
  \item A \textbf{System Security and Privacy Plan (SSP)}: per-system documented security controls
  \item A \textbf{Plan of Action and Milestones (POA\&M)}: tracking of known gaps with remediation plans
\end{itemize}

\begin{center}
\rowcolors{2}{rowgray}{white}
\begin{tabularx}{\textwidth}{L{3.5cm}XL{2.5cm}}
\toprule
\rowcolor{primary}
\tableheader{NIST 800-53 Control Family} & \tableheader{Key Documentation Artifacts Required} & \tableheader{Impact Level} \\
\midrule
Access Control (AC) & Access control policy, user account management procedures, privileged access procedures & Low/Mod/High \\
Audit and Accountability (AU) & Audit policy, log review procedures, retention procedures & Low/Mod/High \\
Risk Assessment (RA) & Risk assessment policy, risk assessment procedures, risk register & Low/Mod/High \\
Incident Response (IR) & IR policy, IR procedures, IR test plan, contact list & Low/Mod/High \\
Configuration Management (CM) & CM policy, baseline configuration procedures, change control procedures & Low/Mod/High \\
System \& Comm. Protection (SC) & Network security policy, encryption standards, boundary protection procedures & Mod/High \\
Program Management (PM) & Information security program policy, SSP procedures, POA\&M management & All \\
\bottomrule
\end{tabularx}
\end{center}

\subsubsection{MCR vs DSR Distinction}

From the Secure Controls Framework:
\begin{itemize}[leftmargin=*]
  \item \textbf{Mandatory Compliance Requirements (MCR):} Externally-influenced; driven by law, regulation, contract (NIST 800-171, CMMC, HIPAA, PCI DSS). These establish the compliance floor.
  \item \textbf{Desired Security Requirements (DSR):} Internally-influenced; based on the organization's risk tolerance and industry context. These are where organizations achieve improved efficiency and security above the floor.
  \item MCR should never be treated as adequate for security --- they are the minimum, not the goal.
\end{itemize}

\subsection{Module 5 Reference: AI-Augmented Documentation}

\subsubsection{Current State of AI in Documentation}

From Google Cloud DORA 2025 State of AI-Assisted Software Development Report:
\begin{itemize}[leftmargin=*]
  \item \textbf{64\%} of software development professionals use AI for writing documentation
  \item \textbf{24.8\%} of developers mostly use AI for creating or maintaining documentation (Stack Overflow 2025)
  \item \textbf{27.3\%} partially use AI for documentation (Stack Overflow 2025)
  \item IBM internal testing: AI-assisted documentation reduced time by \textbf{59\%} on average
  \item 2024 DORA Report: AI-assisted documentation provides \textbf{7.5\% quality improvement}
  \item \textbf{73\%} of developers cite comprehensive API documentation as the primary integration obstacle
\end{itemize}

\subsubsection{AI Documentation Pattern Library}

\begin{center}
\rowcolors{2}{rowgray}{white}
\begin{tabularx}{\textwidth}{L{3.5cm}XL{3cm}}
\toprule
\rowcolor{primary}
\tableheader{Pattern} & \tableheader{Description} & \tableheader{Guardrail Required} \\
\midrule
Inline generation & AI generates docstrings, comments, and inline documentation at commit time & Human review for architectural intent \\
Reference doc synthesis & AI generates API reference from OpenAPI/Swagger specs or code signatures & Accuracy validation against actual behavior \\
SOP first draft & AI generates first-draft SOPs from process descriptions or flowcharts & SME review + walkthrough test \\
Drift detection & AI flags documentation that has not been updated when linked code changes & CI/CD integration; threshold configuration \\
Audience adaptation & AI rewrites technical content for a specified audience persona & Human review for tone and completeness \\
Gap analysis & AI identifies undocumented code paths, endpoints, or procedures & Manual verification of identified gaps \\
\bottomrule
\end{tabularx}
\end{center}

\subsubsection{AI Failure Modes and Mitigations}

\begin{center}
\rowcolors{2}{rowgray}{white}
\begin{tabularx}{\textwidth}{L{4cm}XL{3cm}}
\toprule
\rowcolor{primary}
\tableheader{Failure Mode} & \tableheader{Description \& Risk} & \tableheader{Mitigation} \\
\midrule
``Correct but useless'' generation & Describes what code does without explaining why; lacks architectural intent, edge case rationale & Mandate human review gate; require ``why'' section in templates \\
Missing context & AI generates documentation without understanding system dependencies or deployment constraints & Context injection: provide architecture docs, ADRs, and dependency maps to AI \\
Hallucinated accuracy & AI generates plausible but incorrect behavior descriptions, especially for complex or edge-case logic & Accuracy validation against live system; automated test alignment \\
Documentation debt acceleration & AI generates docs faster than humans can review; review queue grows; approvals become rubber stamps & Enforce review queue limits; block merges when queue exceeds threshold \\
Stale AI confidence & AI-generated docs become authoritative despite drift from code reality & Drift detection metrics in CI pipeline; 90-day staleness flags \\
\bottomrule
\end{tabularx}
\end{center}

\subsection{Safety and Sensitivity Considerations}

\begin{center}
\rowcolors{2}{rowgray}{white}
\begin{tabularx}{\textwidth}{L{4cm}L{3cm}X}
\toprule
\rowcolor{primary}
\tableheader{Area} & \tableheader{Boundary Type} & \tableheader{Guidance} \\
\midrule
Security-sensitive procedures & GATE & Do not publish detailed security procedures (e.g., incident response contact lists, vulnerability handling procedures) in public-facing documentation. Maintain internal-only access controls. \\
PII in documentation examples & ABSOLUTE & Never use real personal data in documentation examples. Use synthetic data generators or named placeholder conventions (e.g., ``Alice Smith,'' ``555-0100''). \\
Access control for compliance docs & GATE & NIST 800-53 and CMMC documentation may contain CUI (Controlled Unclassified Information). Apply appropriate classification and access controls. \\
AI attribution & ANNOTATE & AI-generated documentation sections should be flagged for human review before publication. Do not represent AI output as fully reviewed until it has been. \\
Indigenous / cultural content & ABSOLUTE & If documentation covers Indigenous communities, cultural practices, or traditional knowledge, apply OCAP principles. Nation-specific attribution; sacred knowledge boundaries are absolute gates. \\
\bottomrule
\end{tabularx}
\end{center}

\subsection{Source Bibliography}

\subsubsection{Government \& Agency Sources}

\begin{itemize}[leftmargin=*]
  \item NIST Special Publication 800-53 Revision 5, \textit{Security and Privacy Controls for Information Systems and Organizations}, 2020
  \item NIST Special Publication 800-100, \textit{Information Security Handbook: A Guide for Managers}, 2006
  \item NIST Cybersecurity Framework 2.0, National Institute of Standards and Technology, 2024
  \item NIST Special Publication 800-171 Revision 2, \textit{Protecting Controlled Unclassified Information}, 2021
  \item UK Home Office Engineering Guidance and Standards, \textit{Write Effective Documentation}, 2024 (\url{https://engineering.homeoffice.gov.uk/patterns/write-effective-documentation/})
  \item Utah State University Engineering Writing Center, \textit{Technical Writing Standards}, 2024 (\url{https://engineering.usu.edu/students/ewc/writing-resources/technical-writing-standards})
\end{itemize}

\subsubsection{Industry Standards}

\begin{itemize}[leftmargin=*]
  \item OASIS DITA Technical Committee, \textit{Darwin Information Typing Architecture Version 1.3}, OASIS Standard, 2015
  \item ISO/IEC 27001:2022, \textit{Information Security Management Systems --- Requirements}
  \item ISO/IEC 27002:2022, \textit{Information Security Controls}
  \item Secure Controls Framework (SCF), \textit{Policy vs. Standard vs. Procedure} (\url{https://securecontrolsframework.com})
\end{itemize}

\subsubsection{Research \& Professional Sources}

\begin{itemize}[leftmargin=*]
  \item Google Cloud, \textit{2025 State of AI-Assisted Software Development Report} (DORA Research Program)
  \item Stack Overflow, \textit{2025 Developer Survey}
  \item AltexSoft, \textit{Technical Documentation in Software Development: Types, Best Practices and Tools}, 2024
  \item Paligo, \textit{The Essential Guide to Effective Technical Documentation}, October 2025
  \item Overcast Blog, \textit{AI-Driven Documentation in 2026}, November 2025
  \item Kong Inc., \textit{What is Docs as Code? Guide to Modern Technical Documentation}, April 2025
  \item IBM Think, \textit{AI Code Documentation: Benefits and Top Tips}, March 2026
  \item Secureframe, \textit{NIST 800-53 Required Policies and Procedures}, 2025
  \item Instinctools, \textit{DITA XML: Exploring the Darwin Information Typing Architecture Standard}, May 2024
\end{itemize}

%% ════════════════════════════════════════════════════════════════════════════
%%  STAGE 3 — MISSION PACKAGE
%% ════════════════════════════════════════════════════════════════════════════
\newpage
\stageheader{STAGE 3 --- MISSION PACKAGE}{tertiary}

\section{Milestone Specification}

\subsection{Mission Objective}

Produce a complete, GSD-ready Technical Documentation system for policy and procedure engineering: a taxonomy reference, a structured authoring guide, a docs-as-code pipeline specification, a compliance framework cross-reference, and an AI documentation pattern library --- all assembled as executable GSD skills and reference materials. The output is designed to be consumed by gsd-skill-creator to generate documentation templates, governance models, and pipeline specs for the GSD ecosystem and downstream projects including Fox Infrastructure Group and the College of Knowledge.

\subsection{Architecture Overview}

\begin{asciibox}
MISSION WAVE ARCHITECTURE
==========================

WAVE 0: FOUNDATION (Sequential, Haiku)
  └── Shared taxonomy schema
  └── Source index (URLs, framework refs)
  └── Document templates (policy, procedure, SOP)
  └── Audience persona definitions
         |
         ▼ [CAPCOM GATE 0]
         
WAVE 1A (Parallel, Sonnet)          WAVE 1B (Parallel, Sonnet)
  M1: Taxonomy + Lifecycle           M3: Docs-as-Code Pipeline
  M2: DITA + Structured Authoring    M4: Compliance Frameworks
         |                                   |
         ▼ [CAPCOM GATE 1]                   ▼
         
WAVE 2: SYNTHESIS (Sequential, Opus)
  └── Unified governance model
  └── AI pattern library (M5)
  └── Cross-module integration
  └── GSD skill templates
         |
         ▼ [CAPCOM GATE 2]
         
WAVE 3: PUBLICATION + VERIFICATION (Sequential, Sonnet)
  └── Deliverable assembly
  └── Quality verification matrix
  └── Test execution
  └── Final PDF compilation
\end{asciibox}

\subsection{System Layers}

\begin{enumerate}[leftmargin=*]
  \item \textbf{Foundation Layer:} Shared taxonomy schema, source index, document templates, persona definitions
  \item \textbf{Research Layer:} Five parallel module surveys (M1--M5) producing structured reference content
  \item \textbf{Synthesis Layer:} Cross-module integration producing unified governance model and AI pattern library
  \item \textbf{Publication Layer:} Assembled deliverables, tested against success criteria, formatted and deployed
\end{enumerate}

\subsection{Deliverables}

\begin{center}
\rowcolors{2}{rowgray}{white}
\begin{tabularx}{\textwidth}{C{0.5cm}L{4cm}XL{3cm}}
\toprule
\rowcolor{primary}
\tableheader{\#} & \tableheader{Deliverable} & \tableheader{Success Criteria} & \tableheader{Format} \\
\midrule
1 & Documentation Taxonomy Reference & SC1: complete hierarchy with NIST/ISO/ISACA sourcing & Markdown + table \\
2 & DITA Structured Authoring Guide & SC2: all topic types with decision criteria & Markdown \\
3 & Docs-as-Code Pipeline Spec & SC3: pipeline with tooling, CI/CD templates, drift thresholds & Markdown + YAML \\
4 & Compliance Framework Cross-Reference & SC4: NIST 800-53 control family to artifact mapping & Table (CSV + MD) \\
5 & AI Documentation Pattern Library & SC5: patterns with guardrail specs & Markdown \\
6 & GSD Documentation Skill Template & SC6: policy/procedure/SOP template with structured fields & GSD skill YAML \\
7 & Unified Governance Model & SC7: ownership, review cadence, lifecycle, escalation & Markdown \\
8 & Audience Persona Pack & SC8: 3+ personas with content segmentation guidance & Markdown \\
9 & Research Mission PDF (this document) & SC10: compiles clean, navigable & PDF + .tex \\
10 & Index HTML landing page & Handoff ready & HTML \\
\bottomrule
\end{tabularx}
\end{center}

\subsection{Component Breakdown}

\begin{center}
\rowcolors{2}{rowgray}{white}
\begin{tabularx}{\textwidth}{L{3.5cm}L{3cm}L{2cm}C{2cm}}
\toprule
\rowcolor{primary}
\tableheader{Component} & \tableheader{Dependencies} & \tableheader{Model} & \tableheader{Est. Tokens} \\
\midrule
W0-SCHEMA: Taxonomy schema & None & Haiku & 4k \\
W0-INDEX: Source index & None & Haiku & 2k \\
W0-TEMPLATES: Doc templates & W0-SCHEMA & Haiku & 8k \\
W0-PERSONAS: Audience personas & None & Haiku & 4k \\
W1A-M1: Taxonomy + Lifecycle survey & W0-SCHEMA & Sonnet & 20k \\
W1A-M2: DITA + Structured authoring survey & W0-INDEX & Sonnet & 20k \\
W1B-M3: Docs-as-Code pipeline spec & W0-TEMPLATES & Sonnet & 20k \\
W1B-M4: Compliance framework cross-ref & W0-INDEX, W0-SCHEMA & Sonnet & 25k \\
W2-M5-AI: AI pattern library & W1A-M1, W1B-M3 & Opus & 20k \\
W2-SYNTH: Unified governance model & W1A-M1, W1A-M2, W1B-M3, W1B-M4 & Opus & 30k \\
W2-SKILL: GSD skill template & W2-SYNTH, W0-TEMPLATES & Opus & 15k \\
W3-ASSEMBLE: Deliverable assembly & W2-SYNTH, W2-SKILL, W2-M5-AI & Sonnet & 15k \\
W3-VERIFY: Quality verification & All deliverables & Sonnet & 10k \\
\textbf{TOTAL} & & & \textbf{$\sim$193k} \\
\bottomrule
\end{tabularx}
\end{center}

\subsection{Wave Execution Plan}

\subsubsection{Wave Summary}

\begin{center}
\rowcolors{2}{rowgray}{white}
\begin{tabularx}{\textwidth}{L{2cm}L{2.5cm}L{2.5cm}L{2cm}X}
\toprule
\rowcolor{primary}
\tableheader{Wave} & \tableheader{Mode} & \tableheader{Model} & \tableheader{Duration} & \tableheader{Output} \\
\midrule
Wave 0 & Sequential & Haiku & 1 session & Foundation schemas, templates, personas \\
Wave 1A & Parallel & Sonnet & 1 session & M1 + M2 surveys \\
Wave 1B & Parallel & Sonnet & 1 session & M3 + M4 surveys \\
Wave 2 & Sequential & Opus & 1 session & M5, governance model, GSD skill \\
Wave 3 & Sequential & Sonnet & 1 session & Assembly, verification, publication \\
\bottomrule
\end{tabularx}
\end{center}

\subsubsection{Wave 0: Foundation}

\begin{center}
\rowcolors{2}{rowgray}{white}
\begin{tabularx}{\textwidth}{L{3cm}L{2.5cm}X}
\toprule
\rowcolor{primary}
\tableheader{Task} & \tableheader{Agent} & \tableheader{Specification} \\
\midrule
W0-SCHEMA & Haiku-1 & Generate shared taxonomy schema: JSON structure defining policy/standard/procedure/guideline/control with fields: id, type, owner, lifecycle\_state, review\_cadence, audience, framework\_refs \\
W0-INDEX & Haiku-2 & Build source index: all URLs, framework references, and tool names from Stage 2 bibliography, formatted as structured YAML \\
W0-TEMPLATES & Haiku-3 & Generate three document templates (policy, procedure, SOP) with structured fields: purpose, scope, owner, review\_date, audience, body sections, approval\_block \\
W0-PERSONAS & Haiku-4 & Define five audience personas from Stage 2 Module 1 reference: senior engineer, new hire, compliance auditor, operations staff, product manager \\
\bottomrule
\end{tabularx}
\end{center}

\subsubsection{Wave 1: Parallel Research Tracks}

\textbf{Track A (simultaneous with Track B):}

\begin{center}
\rowcolors{2}{rowgray}{white}
\begin{tabularx}{\textwidth}{L{3cm}L{2.5cm}X}
\toprule
\rowcolor{primary}
\tableheader{Task} & \tableheader{Agent} & \tableheader{Specification} \\
\midrule
W1A-M1 & Sonnet-1 & Produce Module 1 survey: authoritative documentation hierarchy table, document lifecycle state machine, audience persona expansion, technical writing standards checklist. All claims sourced to NIST/ISO/ISACA/USU. \\
W1A-M2 & Sonnet-2 & Produce Module 2 survey: complete DITA topic type reference with examples, DITA vs. Markdown decision matrix, content reuse pattern catalog, DITA-OT pipeline configuration guide. \\
\bottomrule
\end{tabularx}
\end{center}

\textbf{Track B (simultaneous with Track A):}

\begin{center}
\rowcolors{2}{rowgray}{white}
\begin{tabularx}{\textwidth}{L{3cm}L{2.5cm}X}
\toprule
\rowcolor{primary}
\tableheader{Task} & \tableheader{Agent} & \tableheader{Specification} \\
\midrule
W1B-M3 & Sonnet-3 & Produce Module 3 survey: docs-as-code pipeline architecture, complete tooling comparison table, CI/CD YAML job templates (GitHub Actions), drift detection configuration guide, quality metrics dashboard spec. \\
W1B-M4 & Sonnet-4 & Produce Module 4 survey: framework landscape table, NIST 800-53 control family to documentation artifact cross-reference (complete, all 20 families), MCR vs DSR analysis, SSP and POA\&M structure templates. \\
\bottomrule
\end{tabularx}
\end{center}

\subsubsection{Wave 2: Synthesis}

\begin{center}
\rowcolors{2}{rowgray}{white}
\begin{tabularx}{\textwidth}{L{3cm}L{2.5cm}X}
\toprule
\rowcolor{primary}
\tableheader{Task} & \tableheader{Agent} & \tableheader{Specification} \\
\midrule
W2-M5-AI & Opus-1 & Produce Module 5 AI pattern library: six documented patterns with guardrail specifications, five failure mode analyses with mitigations, decision framework for AI vs. human authorship by document type, review gate design. \\
W2-SYNTH & Opus-2 & Synthesize unified governance model integrating M1 taxonomy + M3 pipeline + M4 compliance requirements + M5 AI guardrails. Output: complete governance model document with ownership matrix, review cadence table, lifecycle state machine, and escalation paths. \\
W2-SKILL & Opus-3 & Generate GSD-compatible skill template for documentation generation. Output: YAML skill manifest + Markdown SKILL.md covering: policy template generation, procedure generation, SOP generation, with per-type field specs and validation rules. \\
\bottomrule
\end{tabularx}
\end{center}

\subsubsection{Wave 3: Publication and Verification}

\begin{center}
\rowcolors{2}{rowgray}{white}
\begin{tabularx}{\textwidth}{L{3cm}L{2.5cm}X}
\toprule
\rowcolor{primary}
\tableheader{Task} & \tableheader{Agent} & \tableheader{Specification} \\
\midrule
W3-ASSEMBLE & Sonnet-5 & Assemble all deliverables into organized directory structure. Validate all cross-references. Confirm all 10 deliverables are present and complete. \\
W3-VERIFY & Sonnet-6 & Execute verification matrix: check each success criterion against produced deliverables. Flag any criterion not met. Produce verification report. \\
W3-PUBLISH & Sonnet-7 & Generate index.html landing page. Confirm LaTeX compiles (three-pass xelatex). Produce final package for handoff. \\
\bottomrule
\end{tabularx}
\end{center}

\subsubsection{Cache Optimization Strategy}

\begin{itemize}[leftmargin=*]
  \item Prefix all Wave 1 agent prompts with the full Stage 2 research reference (this document, Stage 2 section) to enable prompt caching
  \item Wave 2 synthesis agents receive Wave 1 outputs as cached context prefix
  \item W0 schemas are injected into all subsequent wave agents as structured YAML (not prose) to minimize token cost
  \item Each agent is a fresh context window --- full relevant context is provided in prompt, not assumed from prior context
\end{itemize}

\subsubsection{Token Budget}

\begin{center}
\rowcolors{2}{rowgray}{white}
\begin{tabularx}{\textwidth}{L{3cm}C{2.5cm}C{2.5cm}C{2.5cm}}
\toprule
\rowcolor{primary}
\tableheader{Wave} & \tableheader{Model} & \tableheader{Est. Tokens} & \tableheader{Context Windows} \\
\midrule
Wave 0 & Haiku & 18k & 1 \\
Wave 1A & Sonnet & 40k & 1 \\
Wave 1B & Sonnet & 45k & 1 \\
Wave 2 & Opus & 65k & 1--2 \\
Wave 3 & Sonnet & 25k & 1 \\
\midrule
\textbf{Total} & \textbf{Mixed} & \textbf{$\sim$193k} & \textbf{5--6} \\
\bottomrule
\end{tabularx}
\end{center}

\subsubsection{Dependency Graph}

\begin{asciibox}
DEPENDENCY GRAPH
================

W0-SCHEMA ──────────────────────────────────┐
W0-INDEX ─────────────────────────────────┐ │
W0-TEMPLATES ───────────────────────────┐ │ │
W0-PERSONAS ────────────────────────────│─│─│─┐
                                        │ │ │ │
                    ┌───────────────────┘ │ │ │
                    │     ┌───────────────┘ │ │
                    ▼     ▼                 │ │
W1A-M1 ──────── (SCHEMA, INDEX) ────────────│─┘
W1A-M2 ──────── (SCHEMA, INDEX)             │
W1B-M3 ──────── (TEMPLATES, INDEX) ─────────┘
W1B-M4 ──────── (SCHEMA, INDEX)

W1A-M1 ────────────────────────────────────┐
W1A-M2 ────────────────────────────────────│──► W2-SYNTH ──► W3-ASSEMBLE
W1B-M3 ────────────────────────────────────│──► W2-M5-AI ──► W3-VERIFY
W1B-M4 ────────────────────────────────────┘──► W2-SKILL ──► W3-PUBLISH

[CAPCOM] ── Gates between: Wave 0 → 1, Wave 1 → 2, Wave 2 → 3
\end{asciibox}

\section{Test Plan}

\subsection{Test Categories}

\begin{center}
\rowcolors{2}{rowgray}{white}
\begin{tabularx}{\textwidth}{L{3cm}C{1.5cm}L{2.5cm}X}
\toprule
\rowcolor{primary}
\tableheader{Category} & \tableheader{Count} & \tableheader{Priority} & \tableheader{On Failure} \\
\midrule
Safety-Critical & 4 & MANDATORY-PASS & BLOCK delivery \\
Core Functionality & 12 & HIGH & Retry or flag for human review \\
Integration & 6 & MEDIUM & Flag; proceed with documented gap \\
Edge Cases & 4 & LOW & Document; address in v1.1 \\
\midrule
\textbf{Total} & \textbf{26} & & \\
\bottomrule
\end{tabularx}
\end{center}

\subsection{Safety-Critical Tests}

\begin{center}
\rowcolors{2}{rowgray}{white}
\begin{tabularx}{\textwidth}{L{1.5cm}L{4cm}X}
\toprule
\rowcolor{primary}
\tableheader{Test ID} & \tableheader{Verifies} & \tableheader{Expected Result / BLOCK Condition} \\
\midrule
SC-01 & No PII in documentation examples & Zero instances of real personal data in any template or example. BLOCK if found. \\
SC-02 & Security-sensitive content access control & Incident response and vulnerability handling procedures must not appear in public-facing deliverables. BLOCK if present in public outputs. \\
SC-03 & Source attribution completeness & Zero unsourced numerical claims. Every statistic must have a traceable citation to a government agency, peer-reviewed source, or professional organization. BLOCK if any unsourced claim found. \\
SC-04 & AI-generated content flagged & Any deliverable section generated primarily by AI (without human-review gate) must be marked ``[PENDING REVIEW]''. BLOCK final publication if unmarked AI content is present. \\
\bottomrule
\end{tabularx}
\end{center}

\subsection{Core Functionality Tests}

\begin{center}
\rowcolors{2}{rowgray}{white}
\begin{tabularx}{\textwidth}{L{1.5cm}L{3cm}L{3cm}X}
\toprule
\rowcolor{primary}
\tableheader{Test ID} & \tableheader{Success Criterion} & \tableheader{Test Method} & \tableheader{Pass Condition} \\
\midrule
CF-01 & SC1: Taxonomy & Check hierarchy table & All 5 terms defined with sourced definitions \\
CF-02 & SC1: Taxonomy & Check examples column & Each term has a concrete, non-generic example \\
CF-03 & SC2: DITA & Check topic types & All 5 DITA topic types documented \\
CF-04 & SC2: DITA & Check decision matrix & DITA vs. Markdown matrix covers 6+ factors \\
CF-05 & SC3: Pipeline & Review pipeline spec & CI/CD YAML template is syntactically valid \\
CF-06 & SC3: Pipeline & Drift detection check & Drift threshold (90 days) is specified and configured \\
CF-07 & SC4: Compliance & Count control families & All 20 NIST 800-53 control families mapped \\
CF-08 & SC4: Compliance & Check artifact columns & Each family has at minimum: policy + procedure listed \\
CF-09 & SC5: AI patterns & Count patterns & Minimum 6 patterns documented \\
CF-10 & SC5: AI patterns & Check guardrails & Each pattern has at least 1 guardrail specification \\
CF-11 & SC6: GSD skill & Validate SKILL.md & Skill file follows GSD skill template format \\
CF-12 & SC7: Governance & Check governance model & Ownership matrix, review cadence, lifecycle, escalation all present \\
\bottomrule
\end{tabularx}
\end{center}

\subsection{Verification Matrix}

\begin{center}
\rowcolors{2}{rowgray}{white}
\begin{tabularx}{\textwidth}{L{1cm}L{5cm}L{3cm}C{2cm}}
\toprule
\rowcolor{primary}
\tableheader{SC\#} & \tableheader{Criterion} & \tableheader{Test IDs} & \tableheader{Status} \\
\midrule
SC1 & Taxonomy clarity & CF-01, CF-02, SC-03 & Pending \\
SC2 & DITA coverage & CF-03, CF-04 & Pending \\
SC3 & Pipeline specification & CF-05, CF-06 & Pending \\
SC4 & Framework mapping & CF-07, CF-08 & Pending \\
SC5 & AI pattern library & CF-09, CF-10, SC-04 & Pending \\
SC6 & GSD skill output & CF-11 & Pending \\
SC7 & Governance model & CF-12 & Pending \\
SC8 & Persona coverage & Integration test & Pending \\
SC9 & Source quality & SC-03 & Pending \\
SC10 & Compilation clean & Build test & Pending \\
SC11 & Test coverage & Count test IDs & Pending \\
SC12 & Handoff ready & Operator review & Pending \\
\bottomrule
\end{tabularx}
\end{center}

\section{Mission Crew Manifest}

\begin{center}
\rowcolors{2}{rowgray}{white}
\begin{tabularx}{\textwidth}{L{3cm}L{3cm}X}
\toprule
\rowcolor{primary}
\tableheader{Role} & \tableheader{Agent (CRAFT)} & \tableheader{Responsibility} \\
\midrule
Mission Commander & CAPCOM (Human) & Wave gate approvals; final quality sign-off \\
Chief Architect & Opus-AGNUS & Synthesis, governance model, AI patterns, GSD skill \\
Module Lead A & Sonnet-DENISE-1 & M1 Taxonomy + Lifecycle survey \\
Module Lead B & Sonnet-DENISE-2 & M2 DITA + Structured authoring survey \\
Module Lead C & Sonnet-DENISE-3 & M3 Docs-as-Code pipeline spec \\
Module Lead D & Sonnet-DENISE-4 & M4 Compliance frameworks survey \\
Assembler & Sonnet-DENISE-5 & Wave 3 deliverable assembly \\
Verifier & Sonnet-DENISE-6 & Quality verification matrix execution \\
Publisher & Sonnet-DENISE-7 & Index page + final packaging \\
Schema Builder & Haiku-PAULA-1 & W0 taxonomy schema \\
Index Builder & Haiku-PAULA-2 & W0 source index \\
Template Builder & Haiku-PAULA-3 & W0 document templates + personas \\
\bottomrule
\end{tabularx}
\end{center}

\section{Execution Summary}

\begin{center}
\rowcolors{2}{rowgray}{white}
\begin{tabularx}{\textwidth}{L{5cm}X}
\toprule
\rowcolor{primary}
\tableheader{Metric} & \tableheader{Value} \\
\midrule
Activation Profile & Squadron (12 roles) \\
Research Modules & 5 \\
Parallel Tracks & 2 (Wave 1A + 1B) \\
CAPCOM Gates & 3 (after W0, after W1, before publication) \\
Total Deliverables & 10 \\
Test Count & 26 (4 safety-critical, 12 core, 6 integration, 4 edge) \\
Estimated Token Budget & $\sim$193k across 5--6 context windows \\
Model Allocation & Opus $\approx$34\%, Sonnet $\approx$56\%, Haiku $\approx$10\% \\
Push Default & Nothing (all commits staged before merge) \\
Bounded Learning & 20\% max change per skill iteration \\
Primary Output Formats & Markdown, YAML, LaTeX PDF, HTML \\
Repository & github.com/Tibsfox/gsd-skill-creator (dev branch) \\
\bottomrule
\end{tabularx}
\end{center}

\vspace{2em}

\begin{quotebox}
\textit{``Documentation is the memory of a system. When it is absent or disconnected from reality, every engineer who follows pays the debt. The spaces between the code are where the meaning lives --- and governing those spaces is as load-bearing as the architecture itself.''}

\vspace{0.8em}
\noindent\textbf{Through-line:} This mission operationalizes the Amiga Principle for knowledge infrastructure: specialized, structured, governed documentation pipelines produce architecturally sound systems. Not brute-force prose generation --- disciplined, testable, versioned artifacts that give people their time back. Documentation treated as engineering. Memory treated as load-bearing structure.
\end{quotebox}

\end{document}
